% =====================================================================
%  Front matter.
%  Source: tier1.md:1-138 (title blurb, status, how-to-study, contents).
%  The <details> "full section list" (tier1.md:73-136) is not dropped --
%  it is reproduced verbatim as Appendix A.
% =====================================================================

\begin{titlepage}
\centering
\vspace*{3.5cm}

{\Huge\bfseries Tier 1 --- Microphysics\par}

\vspace{1.2em}
{\Large Reaction-rate theory, detailed balance, plasma corrections,\\[0.3em]
the weak sector, and the reconciliation of two rate sets\par}

\vspace{0.8em}
{\large\itshape A study report for the conservation-by-construction GNN\\
emulator of silicon burning\par}

\vspace{3em}
\begin{minipage}{0.82\textwidth}
\small
Complete deep dive: where a single number $\lambda_j(T,\rho,\Ye)$ comes from.
Thermonuclear reaction-rate theory from cross-sections up, the REACLIB fit
library and its evaluation (\textbf{S2}), detailed balance and the
partition-function gate (\textbf{S3}), plasma screening (\textbf{S4}), the weak
sector (\textbf{S5}), and the MESA\,/\,pynucastro reconciliation that makes the
whole rate set trustworthy (\textbf{S6}). Everything needed before Tier~2
(stoichiometry, fluxes, and the conservation algebra), in one document.
\end{minipage}

\vfill
{\small Literature audit 2026-08-13\par}
\vspace{0.5em}
{\small Source document: \code{study/tier1.md} \quad$\cdot$\quad
Audit log: \code{study/tier1-audit.md}\par}
\end{titlepage}

% ---------------------------------------------------------------------
\section*{Status of this document}
\addcontentsline{toc}{section}{Status of this document}
\label{sec:status}

\begin{warn}
\textbf{Personal study material, not project spec.} Deliberately written
without reference to \code{docs/} --- everything is derived from first
principles or read off the business code, configs, notebooks, and tests.
\end{warn}

\begin{itemize}
\item Textbook physics and derivations here are \textbf{mine to check}, not
      citable project output. Where I compute a number myself (Gamow peaks,
      screening factors, Fermi energies, coefficient dissections) it is tagged
      \derivedhere{} --- reproducible from the snippets given, but not a
      \code{RESULTS.md} row.

\item Measured project numbers quoted here originate in \code{RESULTS.md} with
      script $+$ commit $+$ data provenance, tagged \resultsref{date}. This
      document is \textbf{not} their source of truth; if a number here
      disagrees with \code{RESULTS.md}, \code{RESULTS.md} wins.

\item \textbf{Literature audit 2026-08-13.} Every scholarly claim, constant,
      and derivation chain was verified against primary sources (arXiv \TeX{}
      cached under \code{data/literature/}, the REACLIB format specification,
      GitHub issues, the local MESA r23.05.1 tree, and pynucastro 2.12.0), with
      all worked arithmetic recomputed in float64; errors found were corrected
      in place. Inline author--year citations point to the References section
      at the end; the claim-by-claim audit log is \code{study/tier1-audit.md}.
\end{itemize}

Notation follows the repo and Tier~0: \nuc{Co}{55} in prose, \code{co55} in
code-adjacent contexts, \Ye{} / $\nu$ / $\varphi$ / $\kappa$ as in the codebase.
Section references of the form \S0.x, \S I.x, \S V.x point back into
\code{study/tier0.md} (and, where that material was itself re-architected, into
the Tier~0 report).

\paragraph{A note on this edition.} The source document is organised
\emph{node by node}: Part~I rate theory, Part~II (S2) REACLIB, Part~III (S3)
detailed balance, Part~IV (S4) screening, Part~V (S5) the weak sector, Part~VI
(S6) reconciliation, Part~VII the guards, Part~VIII the open register. Each of
those parts interleaves four different \emph{kinds} of statement --- a physics
derivation, the code that implements it, the test that pins it, and the policy
that refuses to run without it. That interleaving is the DERIVE\,/\,READ\,/\,%
ORACLE\,/\,SEE study loop the source was written around.

This report re-architects the same material into those four layers:

\begin{center}
\begin{tabularx}{\textwidth}{@{}L{1.5cm} L{3.5cm} Y@{}}
\toprule
\textbf{\S} & \textbf{Layer} & \textbf{Contains} \\
\midrule
\ref{part:ratetheory}--\ref{part:weak} & Theory and derivation &
  the physics of $\lambda_j$: the bare rate, the fitting basis, the reverse
  rate, the plasma corrections, the weak sector \\
\ref{part:evaluator} & The compiled evaluator &
  every code contract, gathered: data structures, the evaluation order, the
  numerical range, the replacements, the tabular path \\
\ref{part:verification} & Verification &
  every oracle, plus the whole MESA\,$\leftrightarrow$\,pynucastro
  reconciliation and how to read its statistics \\
\ref{part:guards} & Executable policy &
  the four guards that refuse to run \\
\ref{part:open} & The open register &
  both audit passes: the physics gaps and the missing Tier~4 \\
\bottomrule
\end{tabularx}
\end{center}

\emph{No content has been omitted.} The cost of layering is that a single
source Part's argument is now split across two or three sections --- Part~III's
algebra lands in \S\ref{part:db}, its spline code in \S\ref{part:evaluator},
its oracles in \S\ref{part:verification}, and its gate in \S\ref{part:guards}.
Three things hold that together: \S\ref{part:evaluator},
\S\ref{part:verification} and \S\ref{part:guards} each open with a table naming
the section that \emph{derived} what they implement, test, or enforce; every
split carries cross-references in both directions; and
Appendix~\ref{app:concordance} is a heading-by-heading concordance from the
source's numbering to this one, including every self-check question that moved
so that no question tests material the reader has not yet reached.
Appendix~\ref{app:sectionmap} reproduces the source document's own section list
verbatim.

Three items were re-homed across parts rather than merely re-layered, because
in the source each was stated somewhere other than where it is derived:

\begin{itemize}
\item \textbf{The two-$\kappa$ conventions rule} (source \S III.10) now sits
      immediately after the screening derivation that produces it
      (\S\ref{sec:two-kappa}).
\item \textbf{Coulomb corrections to NSE} (source \S III.11) sit beside
      screening as the equilibrium-side twin of the same ion-sphere physics
      (\S\ref{sec:nse-coulomb}).
\item \textbf{Screening of the weak sector} (source \S IV.9) moves into the
      weak part (\S\ref{sec:weak-screening}), which removes its forward
      reference to \muE{} entirely.
\end{itemize}

% ---------------------------------------------------------------------
\section*{How to study each node}
\addcontentsline{toc}{section}{How to study each node}
\label{sec:howtostudy}

\begin{figure}[H]
\centering
\begin{tikzpicture}[
  font=\small,
  stepbox/.style={draw, rounded corners=2pt, line width=0.7pt,
               minimum height=8mm, inner xsep=6pt, fill=resultrule!6},
  num/.style={circle, draw, line width=0.7pt, inner sep=1.2pt,
              minimum size=5.2mm, font=\footnotesize\bfseries,
              fill=white}]

\node[stepbox] (a) {\textbf{DERIVE}};
\node[stepbox, below=5mm of a] (b) {\textbf{READ}};
\node[stepbox, below=5mm of b] (c) {\textbf{ORACLE}};
\node[stepbox, below=5mm of c] (d) {\textbf{SEE}};

\begin{scope}[on background layer]
\node[fit=(a)(b)(c)(d), inner sep=0pt] (grp) {};
\end{scope}

\node[num] at ([xshift=-9mm]a.west) {1};
\node[num] at ([xshift=-9mm]b.west) {2};
\node[num] at ([xshift=-9mm]c.west) {3};
\node[num] at ([xshift=-9mm]d.west) {4};

\node[anchor=west, text width=0.66\textwidth] at ([xshift=6mm]a.east)
  {pen-and-paper: get the equation from physics before reading code};
\node[anchor=west, text width=0.66\textwidth] at ([xshift=6mm]b.east)
  {the module, top-to-bottom; its docstring states the contract};
\node[anchor=west, text width=0.66\textwidth] at ([xshift=6mm]c.east)
  {the test file --- it encodes what ``correct'' means, numerically};
\node[anchor=west, text width=0.66\textwidth] at ([xshift=6mm]d.east)
  {the notebook figure --- the measured behaviour on real data};

\draw[-{Latex[length=2mm]}, line width=0.6pt] (a) -- (b);
\draw[-{Latex[length=2mm]}, line width=0.6pt] (b) -- (c);
\draw[-{Latex[length=2mm]}, line width=0.6pt] (c) -- (d);
\end{tikzpicture}
\caption{The four-step study loop for each node of the curriculum. In this
report step~1 lives in \S\ref{part:ratetheory}--\S\ref{part:weak}, step~2 in
\S\ref{part:evaluator}, and step~3 in \S\ref{part:verification}.}
\label{fig:studyloop}
\end{figure}

Tier~1 is the tier where step~1 is heaviest and step~3 is subtlest. Two of its
five nodes (\textbf{S3}, \textbf{S6}) exist \emph{only} because a measurement
contradicted an assumption; their tests encode the contradiction, not a
textbook identity.

% ---------------------------------------------------------------------
\section*{Reading map}
\addcontentsline{toc}{section}{Reading map}
\label{sec:readingmap}

\begin{tabularx}{\textwidth}{@{}L{1.1cm} L{3.7cm} Y@{}}
\toprule
\textbf{\S} & \textbf{Theme} & \textbf{Covers} \\
\midrule
\ref{part:ratetheory} & Rate theory &
  \sigv{} from cross-sections, Coulomb barrier \& Gamow peak, resonances \&
  Hauser--Feshbach, photodisintegration, what the rates \emph{do} in Si
  burning \\
\ref{part:reaclib} & REACLIB (S2) &
  the seven-coefficient form derived, sets\,/\,chapters\,/\,labels, one rate
  dissected, multiple sets and fit validity, a fit library and not a theory \\
\ref{part:db} & Detailed balance (S3) &
  the reverse\,/\,forward ratio derived in full, the phase-space factor per
  $|\Delta N|$, partition functions, the v-flag dissection, $\kappa$ as the DB
  test, gh-575, $Q(T)$ \\
\ref{part:plasma} & Plasma corrections (S4) &
  Debye--H\"uckel derived, $\Gamma$ regimes, \code{chugunov\_2007} term by term,
  the per-reaction pairing and its $\kappa$ consequence, the two-$\kappa$ rule,
  Coulomb terms in NSE \\
\ref{part:weak} & The weak sector (S5) &
  Fermi golden rule $\to$ $ft$ $\to$ allowed rates, Gamow--Teller, EC in a
  degenerate plasma, the $(T,\rhoye)$ tables, weak-sector screening, table
  families, the network asymmetry \\
\ref{part:evaluator} & The compiled evaluator &
  the code contract, the coefficient tensor, the numerical range of $\lambda$,
  the v-flag replacement, the tabular path, the ledger enforcement \\
\ref{part:verification} & Verification (S6) &
  the oracles, canonical keys, membership vs values, dispositions, the measured
  comparison, the Fortran probe, what it licenses \\
\ref{part:guards} & Executable policy &
  \code{fluxes/guards.py}: the four Tier-1 gates as code that refuses to run \\
\ref{part:open} & Open prerequisites &
  standing cross-tier register; two audit passes: the physics the code
  implements (ten gaps) and the project's own spec documents vs the study map
  (an entire missing \textbf{Tier 4}) \\
\bottomrule
\end{tabularx}

\vspace{1em}
{\small The source document's own contents table and full section list --- its
original Part~I--VIII architecture --- are reproduced verbatim in
Appendix~\ref{app:sectionmap}.}

\clearpage
\tableofcontents
\clearpage
